% ---------------------------------------------------------------------------
%  figs/tikz_figures.tex  — all standalone TikZ figures for paperopenmp2
% ---------------------------------------------------------------------------

% ============================================================
% FIGURE 1 — Sequential BFS zone-graph exploration
% ============================================================
\begin{figure}[t]
\centering
\begin{tikzpicture}[node distance=10mm and 14mm]

  % ---- Waiting queue (left column) ---------------------------------
  \node[cpubox, minimum width=34mm] (w1) {$(\ell_1,\, Z_a)$};
  \node[cpubox, minimum width=34mm, below=4mm of w1] (w2) {$(\ell_2,\, Z_b)$};
  \node[cpubox, minimum width=34mm, below=4mm of w2] (w3) {$(\ell_1,\, Z_c)$};
  \node[below=2mm of w3, font=\scriptsize\itshape, text=black!50] (wdots) {$\vdots$};
  \node[above=1mm of w1, font=\small\bfseries, text=cA] {$\Wait$ (deque)};

  \draw[decorate, decoration={brace,amplitude=4pt}, thick, draw=black!40]
    ([xshift=-2mm]w1.north west) -- ([xshift=-2mm]wdots.south west)
    node[midway, left=4mm, font=\scriptsize\itshape, text=black!50] {FIFO};

  % ---- Processing box (centre) ------------------------------------
  \node[box, minimum width=38mm, minimum height=14mm,
        right=20mm of w2, fill=cpu!8, draw=cpu!60] (proc)
    {expand\\$(\ell,Z) \to \succfn(\ell,Z)$};

  % ---- Passed list (right column) ----------------------------------
  \node[mergebox, minimum width=34mm, right=20mm of proc, yshift=8mm] (p1)
    {$\Passed[\ell_1] = \{Z_a, Z_c, \ldots\}$};
  \node[mergebox, minimum width=34mm, below=4mm of p1] (p2)
    {$\Passed[\ell_2] = \{Z_b, \ldots\}$};
  \node[below=2mm of p2, font=\scriptsize\itshape, text=black!50] {$\vdots$};
  \node[above=1mm of p1, font=\small\bfseries, text=cC] {$\Passed$ (per-location)};

  % arrows
  \draw[flow] (w2.east) -- node[above, lbl]{dequeue} (proc.west);
  \draw[flow] (proc.east) -- node[above, lbl]{subsumption\\check} (p2.west);
  \draw[dflow] (p2.west) to[bend right=12] node[below, lbl]{not subsumed:\\enqueue} (w3.east);

  % target check
  \node[cpubox, minimum width=28mm, below=18mm of proc] (tgt)
    {target check\\$\ell=\ell^\star \wedge Z\cap\phi\neq\emptyset$?};
  \draw[flow] (proc.south) -- node[right, lbl]{new successor} (tgt.north);
  \node[right=8mm of tgt, font=\small\bfseries, text=red!70!black] (yes) {\textbf{SAT}};
  \draw[flow, draw=red!60] (tgt.east) -- node[above,lbl]{yes} (yes);

\end{tikzpicture}
\caption{Sequential BFS zone-graph exploration.  Each symbolic state $(\ell,Z)$ is
dequeued, its successors computed, and each successor either discarded (subsumed by
an existing zone in $\Passed[\ell']$) or enqueued for further expansion.}
\label{fig:seq-bfs}
\end{figure}

% ============================================================
% FIGURE 2 — Level-synchronous parallel BFS (OpenMP)
% ============================================================
\begin{figure}[t]
\centering
\begin{tikzpicture}[node distance=8mm and 12mm]

  % ---- Frontier box ------------------------------------------------
  \node[cpubox, minimum width=56mm, minimum height=9mm] (front)
    {$\Frontier_k = \bigl\{(\ell_i, Z_i)\bigr\}_{i=1}^{N}$};
  \node[above=1mm of front, font=\small\bfseries, text=cA]
    {Frontier level $k$ (shared read-only array)};

  % ---- Thread boxes ------------------------------------------------
  \node[threadbox, below=14mm of front, xshift=-42mm] (t0) {thread 0\\$\mathit{buf}_0$};
  \node[threadbox, below=14mm of front, xshift=-14mm] (t1) {thread 1\\$\mathit{buf}_1$};
  \node[threadbox, below=14mm of front, xshift=14mm]  (t2) {thread 2\\$\mathit{buf}_2$};
  \node[threadbox, below=14mm of front, xshift=42mm]  (t3) {$\cdots$\\$\mathit{buf}_{T-1}$};

  % arrows frontier -> threads
  \draw[flow, draw=cA!70] (front.south) -- ++(0,-4mm) -| (t0.north);
  \draw[flow, draw=cA!70] (front.south) -- ++(0,-4mm) -| (t1.north);
  \draw[flow, draw=cA!70] (front.south) -- ++(0,-4mm) -| (t2.north);
  \draw[flow, draw=cA!70] (front.south) -- ++(0,-4mm) -| (t3.north);

  % omp annotation
  \node[right=4mm of front.south east, font=\scriptsize\ttfamily, text=cpu,
        yshift=-7mm, anchor=west]
    {\#pragma omp parallel for schedule(dynamic,8)};

  % ---- Merge / filter box ------------------------------------------
  \node[mergebox, minimum width=56mm, minimum height=9mm,
        below=14mm of t1, xshift=14mm] (merge)
    {sequential merge: subsumption filter + $\Passed$ update};
  \node[below=0.5mm of merge.south, text=cGrey, font=\scriptsize\itshape]
    {(single-threaded critical section)};

  % arrows threads -> merge
  \draw[flow, draw=cC!70] (t0.south) |- (merge.west);
  \draw[flow, draw=cC!70] (t1.south) -- (t1.south |- merge.north);
  \draw[flow, draw=cC!70] (t2.south) -- (t2.south |- merge.north);
  \draw[flow, draw=cC!70] (t3.south) |- (merge.east);

  % ---- Next frontier -----------------------------------------------
  \node[cpubox, minimum width=56mm, minimum height=9mm,
        below=12mm of merge] (next)
    {$\Frontier_{k+1}$ (new unseen states)};
  \draw[flow] (merge.south) -- node[right,lbl]{not subsumed} (next.north);

  % loop-back arrow
  \draw[dflow, draw=cA!50]
    (next.east) -- ++(8mm,0) |- node[right, lbl, yshift=12mm]{repeat until $\Frontier_{k+1}=\emptyset$}
    (front.east);

  % LU extrapolation note
  \node[box, fill=cB!8, draw=cB!50, minimum width=40mm,
        right=18mm of merge] (ext)
    {$\extraLU$ applied inside\\each $\succfn()$ call};
  \draw[dflow, draw=cB!50] (merge.east) -- (ext.west);

\end{tikzpicture}
\caption{Level-synchronous parallel BFS.  All states in $\Frontier_k$ are expanded
in parallel by $T$ OpenMP threads, each writing successors into a private buffer
$\mathit{buf}_t$.  The buffers are merged sequentially: subsumption against $\Passed$
is checked and new states that survive form $\Frontier_{k+1}$.}
\label{fig:par-bfs}
\end{figure}

% ============================================================
% FIGURE 3 — DBM representation of a zone
% ============================================================
\begin{figure}[t]
\centering
\begin{tikzpicture}

  \matrix (M) [matrix of nodes,
               nodes={draw=black!30, minimum width=10mm, minimum height=8mm,
                      font=\small, inner sep=2pt, anchor=center},
               column sep=0pt, row sep=0pt,
               row 1/.style={nodes={fill=cpu!15, font=\small\bfseries}},
               column 1/.style={nodes={fill=cpu!15, font=\small\bfseries}}]
  {
    $-$     & $x_0$ & $x_1$ & $x_2$ \\
    $x_0$   & $0$   & ${\leq}b_{01}$ & ${\leq}b_{02}$ \\
    $x_1$   & ${\leq}b_{10}$ & $0$ & ${\leq}b_{12}$ \\
    $x_2$   & ${\leq}b_{20}$ & ${\leq}b_{21}$ & $0$ \\
  };

  \node[right=8mm of M, align=left, font=\small] (ann) {
    $b_{ij} \in \mathbb{Z} \cup \{+\infty\}$\\[3pt]
    each entry encodes\\
    $x_i - x_j \leq b_{ij}$\\[4pt]
    \textit{canonical form:}\\
    \textit{shortest paths}
  };

  \draw[draw=cC, thick, rounded corners=1pt]
    (M-2-2.north west) rectangle (M-2-4.south east);
  \node[below=2mm of M.south, font=\scriptsize\itshape, text=cC]
    {row $x_0$: upper bounds on individual clocks};

\end{tikzpicture}
\caption{A Difference Bound Matrix (DBM) for a $2$-clock automaton.  Clock $x_0 \equiv 0$
is the reference; entry $b_{ij}$ encodes $x_i - x_j \leq b_{ij}$.  Canonical form is
maintained after every operation so zone inclusion runs in $O(n^2)$.}
\label{fig:dbm}
\end{figure}

% ============================================================
% FIGURE 4 — Diagonal splitting overview
% ============================================================
\begin{figure}[t]
\centering
\begin{tikzpicture}[node distance=10mm and 18mm]

  \begin{scope}
    \node[statebox] (l0) {$\ell_0$};
    \node[statebox, right=22mm of l0] (l1) {$\ell_1$};
    \node[statebox, right=22mm of l1] (l2) {$\ell_2$};
    \draw[flow] (l0) -- node[above, font=\scriptsize]{$x_1 - x_2 \leq 3,\; x_1{:=}0$} (l1);
    \draw[flow] (l1) -- node[above, font=\scriptsize]{$x_2 \leq 5$} (l2);
    \node[above=5mm of l1, font=\small\bfseries] {(a) Original TA with diagonal guard};
  \end{scope}

  \begin{scope}[yshift=-38mm]
    \node[statebox] (s00) {$(\ell_0,0)$};
    \node[statebox, right=24mm of s00] (s10) {$(\ell_1,0)$};
    \node[statebox, right=24mm of s10] (s11) {$(\ell_1,1)$};
    \node[statebox, right=18mm of s11] (s2)  {$(\ell_2,{*})$};

    \draw[flow] (s00) -- node[above, font=\scriptsize]
      {$d_1{=}0$: $-x_2\leq 3$,\; $x_1{:=}0$} (s10);
    \draw[flow] (s00) to[bend left=16] node[above, font=\scriptsize]
      {$d_1{=}1$: $-x_2\leq 3$,\; $x_1{:=}0$} (s11);
    \draw[flow] (s10) -- node[above, font=\scriptsize]{$x_2\leq 5$} (s2);
    \draw[flow] (s11) to[bend left=10] node[below, font=\scriptsize]{$x_2\leq 5$} (s2);

    \node[above=5mm of s10, font=\small\bfseries] {(b) After diagonal splitting};
    \node[below=4mm of s10, font=\scriptsize\itshape, text=black!55]
      {$d_1 = [x_1 - x_2 \leq 3]$};
  \end{scope}

  \draw[dashed, black!30] (-1,-18mm) -- (13,-18mm);

\end{tikzpicture}
\caption{Diagonal splitting.  The diagonal guard $x_1 - x_2 \leq 3$ is encoded as a
Boolean bit $d_1$ in the extended location.  After resetting $x_1$, the diagonal
becomes a non-diagonal bound $-x_2 \leq 3$, yielding a diagonal-free automaton on
which $\extraLU$ is sound.}
\label{fig:diag-split}
\end{figure}

% ============================================================
% FIGURE 5 — Thread-local buffer layout (memory view)
% ============================================================
\begin{figure}[t]
\centering
\newcommand*{\threadww}{18mm}%
\newcommand*{\threadhh}{24mm}%
\begin{tikzpicture}[node distance=6mm and 10mm]

  \node[cpubox, minimum width=88mm, minimum height=10mm] (front)
    {Shared frontier $\Frontier_k[0\ldots N{-}1]$\quad\textit{(read-only during parallel phase)}};
  \node[above=1mm of front, font=\small\bfseries, text=cA]{Frontier level $k$};

  % thread 0
  \node[rounded corners=2pt, draw=cA!50, fill=cpu!20, thick,
        minimum width=\threadww, minimum height=\threadhh, align=center,
        below=14mm of front, xshift=-33mm, font=\scriptsize] (tha) {};
  \node[font=\scriptsize\bfseries, text=cA, anchor=north] at (tha.north) [yshift=-2mm] {thread 0};
  \node[zonenode] at ($(tha.north)+(0,-8mm)$)  {$(\ell, Z)$};
  \node[zonenode] at ($(tha.north)+(0,-14mm)$) {$(\ell',Z')$};
  \node[font=\tiny, text=black!40] at ($(tha.north)+(0,-18mm)$) {$\vdots$};
  \node[font=\tiny\itshape, text=black!55, anchor=south] at (tha.south) [yshift=1mm] {$\mathit{buf}_0$};

  % thread 1
  \node[rounded corners=2pt, draw=cA!50, fill=cpu!12, thick,
        minimum width=\threadww, minimum height=\threadhh, align=center,
        below=14mm of front, xshift=-11mm, font=\scriptsize] (thb) {};
  \node[font=\scriptsize\bfseries, text=cA, anchor=north] at (thb.north) [yshift=-2mm] {thread 1};
  \node[zonenode] at ($(thb.north)+(0,-8mm)$)  {$(\ell_2,Z_a)$};
  \node[zonenode] at ($(thb.north)+(0,-14mm)$) {$(\ell_3,Z_b)$};
  \node[font=\tiny, text=black!40] at ($(thb.north)+(0,-18mm)$) {$\vdots$};
  \node[font=\tiny\itshape, text=black!55, anchor=south] at (thb.south) [yshift=1mm] {$\mathit{buf}_1$};

  % thread 2
  \node[rounded corners=2pt, draw=cA!50, fill=cpu!06, thick,
        minimum width=\threadww, minimum height=\threadhh, align=center,
        below=14mm of front, xshift=11mm, font=\scriptsize] (thc) {};
  \node[font=\scriptsize\bfseries, text=cA, anchor=north] at (thc.north) [yshift=-2mm] {thread 2};
  \node[zonenode] at ($(thc.north)+(0,-8mm)$)  {$(\ell_1,Z_c)$};
  \node[font=\tiny, text=black!40] at ($(thc.north)+(0,-12mm)$) {$\vdots$};
  \node[font=\tiny\itshape, text=black!55, anchor=south] at (thc.south) [yshift=1mm] {$\mathit{buf}_2$};

  % thread T-1
  \node[rounded corners=2pt, draw=cA!50, fill=cGrey!12, thick,
        minimum width=\threadww, minimum height=\threadhh, align=center,
        below=14mm of front, xshift=33mm, font=\scriptsize] (thd) {};
  \node[font=\scriptsize\bfseries, text=cA, anchor=north] at (thd.north) [yshift=-2mm] {thread $T{-}1$};
  \node[font=\tiny, text=black!40] at ($(thd.north)+(0,-10mm)$) {$\vdots$};
  \node[font=\tiny\itshape, text=black!55, anchor=south] at (thd.south) [yshift=1mm] {$\mathit{buf}_{T-1}$};

  % arrows frontier -> threads
  \draw[flow, draw=cA!70] (front.south) -- ++(0,-4mm) -| (tha.north);
  \draw[flow, draw=cA!70] (front.south) -- ++(0,-4mm) -| (thb.north);
  \draw[flow, draw=cA!70] (front.south) -- ++(0,-4mm) -| (thc.north);
  \draw[flow, draw=cA!70] (front.south) -- ++(0,-4mm) -| (thd.north);

  % merge
  \node[mergebox, minimum width=88mm, minimum height=12mm,
        below=10mm of thb, xshift=11mm] (merge)
    {sequential merge:\quad
     $s \notin \Passed \;\Rightarrow\; \Passed \mathrel{+}= s,\;
      \Frontier_{k+1} \mathrel{+}= s$};

  \draw[flow, draw=cC!70] (tha.south) |- (merge.west);
  \draw[flow, draw=cC!70] (thb.south) -- (thb.south |- merge.north);
  \draw[flow, draw=cC!70] (thc.south) -- (thc.south |- merge.north);
  \draw[flow, draw=cC!70] (thd.south) |- (merge.east);

\end{tikzpicture}
\caption{Memory layout during one parallel frontier expansion.  Each thread $t$ holds a
private buffer into which it appends the successors of the states it owns.  No
synchronisation is needed during expansion; all shared-state updates happen in the
single-threaded merge phase.}
\label{fig:mem-layout}
\end{figure}
